\subsection{Background}

Voice-based virtual assistants have advanced rapidly alongside broader progress in Artificial Intelligence, evolving from simple command-response systems into sophisticated conversational agents capable of complex reasoning and task execution. Despite their growing capabilities, widespread adoption has surfaced critical concerns regarding user privacy and security.

In particular, sensitive user information such as behavioral patterns, personal preferences, and private data must be safeguarded against unauthorized access. Speaker verification (SV) determines whether a query voice matches a claimed identity and is therefore suitable for authenticating sensitive requests. Speaker identification (SID) determines which enrolled speaker is talking and supports personalization.

Beyond security, robust Speaker Verification also enables richer, more personalized user experiences. By recognizing individual speakers, a virtual assistant can tailor its responses, recommendations, and guidance to each user's unique context and history.

However, Speaker Identification in real-world settings remains a challenging task. Spoken utterances vary considerably due to factors such as background noise, recording conditions, emotional state, and dialectal differences. Moreover, models trained on data from one linguistic or acoustic domain may generalize poorly to others, underscoring the need for evaluation across diverse and representative datasets. These challenges motivate the adoption of modern deep learning architectures that can learn robust speaker representations directly from raw or lightly processed audio.

Motivated by these considerations, this project provides a concise introduction to Speaker Identification and presents a systematic analysis and comparison of several methods and models grounded in Statistical Machine Learning, with the overarching goal of advancing Speaker Verification for reliable and efficient virtual assistants.

\subsection{Our Contributions}
This work provides three key contributions:

\begin{enumerate}
    \item \textbf{Dataset curation.} We benchmark speaker verification across two complementary corpora: the regional \emph{ViMD} dataset, capturing Vietnamese multi-dialectal variations, and the multilingual \emph{TidyVoice2026} corpus, assessing cross-lingual robustness under standardized verification protocols.

    \item \textbf{Model benchmarking.} We benchmark two speaker embedding architectures---RawNet3 and ECAPA-TDNN---after target-domain fine-tuning on both datasets. All reported EER, minDCF, FAR, FRR, and TAR results are obtained under a shared speaker-verification protocol.

    \item \textbf{Proof-of-concept application.} We develop \emph{Who-Speak-AI}, a working voice-based virtual assistant that demonstrates secure, speaker-aware interaction. The application adopts a modular pipeline comprising Voice Activity Detection (VAD), Automatic Speech Recognition (ASR), a Large Language Model (LLM) for dialogue, and Text-to-Speech (TTS) synthesis. This end-to-end prototype illustrates how Speaker Verification can be seamlessly integrated into a practical conversational system.
\end{enumerate}

\subsection{Report Organization}
The remainder of this report is organized as follows.
\textbf{Section~\ref{sec:related}} reviews related works on Speaker Verification and recognition.
\textbf{Section~\ref{sec:datasets}} describes the two datasets used in our experiments.
\textbf{Section~\ref{sec:model}} presents the architectures of the models under study.
\textbf{Section~\ref{sec:methods}} details the experimental setup and evaluation methodology.
\textbf{Section~\ref{sec:results}} reports and analyzes the experimental results.
\textbf{Section~\ref{sec:application}} introduces the Who-Speak-AI application as a proof-of-concept deployment.
Finally, \textbf{Section~\ref{sec:conclusion}} concludes the report and discusses future directions.
